\begin{textsample}{NEG Dim 8 – human – Score -1.00 – sp/sp\_inf\_e\_ef\_intro\_9\_p5.txt}  \label{ex:f8_neg_012}
A avaliação integra e constitui um espaço crítico-reflexivo da prática docente.
Deve garantir coerência com os princípios pedagógicos que orientam o desenvolvimento pleno dos estudantes.
No processo avaliativo, é necessário que se considerem as aprendizagens prescritas no Currículo Paulista.
Na Educação Infantil, a avaliação deve ser realizada por meio de observações e dos mais diversos registros, conforme a Lei de Diretrizes e Bases da Educação, na seção 11, artigo 31, que diz que “ [ … ] a avaliação far -se-á mediante o acompanhamento e registro do seu desenvolvimento, sem o objetivo de promoção, mesmo para o acesso ao Ensino Fundamental ”.
Como \textbf{exemplo} de registros, podemos citar : relatórios, fotografias, filmagens, produções infantis, diários, portfólios, murais, entre outros.

% matched lemmas: exemplo
\end{textsample}
